\section{总结与延伸}

\subsection{全讲总结表}
\begin{table}[H]
\centering
\begin{tabular}{p{3.0cm}p{4.6cm}p{4.0cm}}
\toprule
\textbf{主题} & \textbf{核心观点} & \textbf{工程动作} \\
\midrule
问题定义 & Agent 目标从答题转向真实任务闭环 & 建立环境交互评测而非仅离线评分 \\
Benchmark & OS-World 补齐真实计算机场景缺口 & 任务脚本化、环境可重放、结果可验证 \\
数据体系 & 轨迹数据是主要瓶颈 & 教程驱动采集 + 合成增强 + 质量治理 \\
模型体系 & 统一感知/规划/执行并分阶段训练 & 引入规划信号，强化 grounding 与长程决策 \\
落地策略 & 生产可用性依赖可观测与可恢复 & 增加审计、回滚、权限和人工接管机制 \\
\bottomrule
\end{tabular}
\caption{课程结论的工程化抽象}
\end{table}

\subsection{进一步阅读}
\begin{itemize}
    \item OS-World 论文与项目文档（真实计算机环境 Agent benchmark）
    \item WebArena / VisualWebArena / Mind2Web（Web Agent 评测演进）
    \item GUI grounding 与 action modeling 相关工作（视觉定位到动作执行）
    \item 多模态 Agent 数据合成与轨迹筛选论文（可扩展训练数据管线）
    \item Agent 安全与审计实践（权限、回滚、可观测性）
\end{itemize}

\subsection{延伸问题}
\begin{itemize}
    \item 如何设计跨应用任务的统一中间表示，减少平台耦合？
    \item 在在线生产环境中，如何平衡自动化效率与安全边界？
    \item 当任务目标不可完全脚本化时，评估系统如何与人工反馈协同？
\end{itemize}

\end{document}
